\section{Weyl 代数$A_n$与 $A_n$-模}
Weyl代数$A_n$是指仿射空间$\left(\mathbb{A}^n_k,\mathcal{O}\right)$上的微分算子层（$k$为特征零域），是$\mathcal{D}$-模理论中的最简单、基本对象。许多定理命题的证明以及具体的计算都要现在Weyl代数$A_n$上进行。因此在本节中我们将系统地讨论Weyl代数$A_n$的定义、基本性质以及$A_n$-模的结构，为后续章节中$\mathcal{D}$-模理论中的特征簇、$b$-函数等问题的研究奠定基础。


\subsection{定义和基本性质}
设 $k[X]=k[x_1,\dots,x_n]$ 是特征为零的域 $k$ 上
$n$ 个可交换不定元的多项式环，
记其所有 $k$-线性自同态所成的代数为 $\End_k(k[X])$.
对每个 $i=1,\dots,n$,
定义算子 $\hat{x}_i\in \End_k(k[X])$ 为
\[
    \hat{x}_i(f)=x_i f, \qquad f\in k[X].
\]
同时定义偏导算子
\[
    \partial_i(f)=\frac{\partial f}{\partial x_i}.
\]
\begin{definition}
    由算子 $\hat{x}_1,\dots,\hat{x}_n$ 与
    $\partial_1,\dots,\partial_n$
    生成的 $\End_k(k[X])$的$k$-子代数
    称为\emph{$n$阶Weyl代数},
    记为 $A_n(k)$.
    约定 $A_0(k)=k$.
\end{definition}



\begin{proposition}
    在 $A_n(k)$ 中有交换关系
    \[
        [\partial_i,\hat{x}_j]=\delta_{ij}, \qquad
    [\partial_i,\partial_j]=[\hat{x}_i,\hat{x}_j]=0,
    \]
    并且对任意 $f\in k[X]$,
    \[
        [\partial_i,f]=\frac{\partial f}{\partial x_i}.
    \]
\end{proposition}

\begin{proposition}\label{prop:weyl-basis}
    给定域$k$，与$k$上的Weyl代数$A_n(k)$。
    \begin{enumerate}
        \item 对任意多重指标 $\alpha,\beta\in\mathbb{N}^n$,
              有
              \[
                  \partial^\beta x^\alpha
                  =
                  \sum_{\beta_1+\beta_2=\beta}
                  \binom{\beta}{\beta_1}
                  \frac{\partial x^\alpha}{\partial x^{\beta_1}}
                  \partial^{\beta_2}.
              \]
              因此
              \[
                  \partial^\beta x^\alpha
                  =
                  x^\alpha\partial^\beta+D',
                  \quad
                  \deg D'\le |\alpha|+|\beta|-2.
              \]

        \item 集合
              \[
                  \{x^\alpha\partial^\beta:
                  \alpha,\beta\in\mathbb{N}^n\}
              \]
              构成 $A_n(k)$ 作为 $k$-向量空间的一组基。
    \end{enumerate}
    \begin{proof}
        （1）即为 Leibniz 公式的特殊情形，直接验证即可。

        （2）注意到$A_n(k)$中元素的交换性，$A_n(k)$中的任意元素都可以表示为有限和
        \begin{equation*}
            \sum c x_1^{\alpha_1}\partial_1^{\beta_1}\cdots x_n^{\alpha_n}\partial_n^{\beta_n}
        \end{equation*}
        对每个分量$i$，利用（1）中的公式可以将$x_i^{\alpha_i}\partial_i^{\beta_i}$表示$\left\{x_i^s\partial_i^t\right\}$的线性组合。代入上式 中，得到$x^\alpha\partial^\beta$的线性组合。线性无关性容易验证，因此$\{x^\alpha\partial^\beta\}$是$A_n(k)$的一组基。
    \end{proof}
\end{proposition}

Weyl代数$A_n(k)$的结构可以通过自由代数的商来描述。

\begin{theorem}
    \label{thm:A}
    设 $k\langle z_1,\dots,z_{2n}\rangle$
    为自由结合代数，
    $J$ 为由关系
    \[
        [z_{i+n},z_i]-1,\quad
        [z_i,z_j]\quad (j\neq i+n)
    \]
    生成的双边理想，则存在 $k$-代数同构
    \[
        A_n(k)\cong k\langle z_1,\dots,z_{2n}\rangle/J.
    \]
\end{theorem}

\begin{proof}
    定义满同态
    $\phi(z_i)=x_i,\ \phi(z_{i+n})=\partial_i$.
    由 Weyl 关系知 $J\subset\ker\phi$,
    从而诱导同态
    $\bar\phi$.
    由命题 \ref{prop:weyl-basis},
    标准单项式在商代数中线性无关，
    故 $\bar\phi$ 为同构。
\end{proof}

\subsection{次数与阶}

\begin{definition}
    由\ref{prop:weyl-basis}，每个元素 $D\in A_n$ 都可以唯一表示为
    \begin{equation*}
        D=\sum_{\alpha,\beta}
        a_{\alpha\beta}x^\alpha\partial^\beta\in A_n.
    \end{equation*}
    称为$D$的\emph{典范形式}。
    定义 $D$ 的\emph{次数}为
    \[
        \deg(D)=\max\{|\alpha|+|\beta|:
        a_{\alpha\beta}\neq0\},
    \]
    并约定零元的次数为 $-\infty$.
\end{definition}

\begin{theorem}
    对任意 $D,D'\in A_n$，有：
    \begin{enumerate}
        \item
              \[
                  \deg(D+D')\le\max\{\deg D,\deg D'\},
              \]
              且当 $\deg D\neq\deg D'$ 时取等号；
        \item
              \[
                  \deg(DD')=\deg D+\deg D'；
              \]
        \item
              \[
                  \deg([D,D'])\le\deg D+\deg D'-2.
              \]
    \end{enumerate}
    \begin{proof}
        利用\ref{prop:weyl-basis}和定义，直接验证即可。
    \end{proof}
\end{theorem}

\begin{corollary}
    Weyl 代数 $A_n$ 是整环。
\end{corollary}

\begin{definition}
    若考虑典范形式
    \[
        D=\sum_{\alpha,\beta}
        a_{\alpha\beta}x^\alpha\partial^\beta,
    \]
    则定义 $D$ 的\emph{阶}为
    \[
        \ord(D)=\max\{|\beta|:a_{\alpha\beta}\neq0\}.
    \]
\end{definition}
与次数函数类似，阶函数也满足一些基本性质。
\begin{proposition}
    阶函数满足：
    \begin{enumerate}
        \item
              $\ord(P+Q)\le\max\{\ord P,\ord Q\}$；
        \item
              $\ord(PQ)=\ord P+\ord Q$；
        \item
              $\ord([P,Q])\le\ord P+\ord Q-1$.
    \end{enumerate}
\end{proposition}

利用$A_n$上的度数函数，我们可以证明$A_n$的单性。

\begin{theorem}
    \label{thm:simple}
    Weyl 代数 $A_n$ 是单代数。
\end{theorem}
\section{Jacobian 猜想}
给定特征为零的域$k$和正整数$n$。设仿射空间$\mathbb{A}_n$上的多项式映射
\begin{equation*}
    f:\mathbb{A}_n\to \mathbb{A}_n
\end{equation*}
用$J(f)$表示其Jacobian矩阵。如果$f$可逆（在多项式映射意义下，下述同理），则$\Delta f:=\det J(f)$必须是一个非零常数。
Jacobian猜想断言，反过来，如果$\Delta f$是一个非零常数（不失一般性地假设$\Delta f=1$)，那么$f$必定是可逆的。
\begin{example}
    设
    \begin{equation*}
        f(x_1,x_2)=(x_1+x_2^2,x_2)
    \end{equation*}
    ，则$\Delta f=1$，且$f^{-1}(y_1,y_2)=(y_1-y_2^2,y_2)$.
\end{example}

通过下述命题\ref{lem:Bernstein-inequality}和\ref{lem:Jacobian}，我们可以将Jacobian猜想等价地表述为一个的代数问题。
\begin{proposition}
    \label{pro:Jacobian}
    多项式映射$f:\mathbb{A}_n\to \mathbb{A}_n$，当且仅当其余态射$f^\sharp$为$\mathcal{O}_{\mathbb{A}_n}=k[x_1,\dots,x_n]$上的环同构。
\end{proposition}



\begin{lemma}
    \label{lem:Jacobian}
    设 $f: \mathbb{A}_n \to \mathbb{A}_n$ 为多项式映射，$f_1,\dots,f_n$ 为其分量函数，且$\Delta f\neq 0$   处处成立。
    则余态射 $f^\sharp$ 为单射。
\end{lemma}

\begin{proof}
    假设引理不成立，则存在最低次数非零多项式
    $\varphi$
    满足

    \begin{equation*}
        f^\sharp(\varphi)=\varphi\circ f=0
    \end{equation*}
    由链式法则得
    \begin{equation*}
        \nabla (\varphi\circ f)
        =
        (\nabla \varphi)(f)\cdot J(f)=0。
    \end{equation*}
    由于 $J(f)$可逆，
    推出
    \begin{equation*}
        \nabla (\varphi\circ f)=0。
    \end{equation*}
    但存在某个$i$ 使得 $\frac{\partial \varphi}{\partial x_i}$ 非零，且次数严格小于 $\varphi$，与极小性矛盾。
\end{proof}
考虑\ref{lem:Jacobian}中$f^{\sharp}$的像$k[F_1,\dots,F_n]\subset k[x_1,\dots,x_n]$，我们可以将Jacobian猜想等价地表述为：当$\Delta f=1$时，$k[F_1,\dots,F_n]=k[x_1,\dots,x_n]$。

\subsection{Jacobian猜想在特定情况下的证明}

\begin{definition}
    设 $D$ 为 $k$-代数 $A$ 的导子。
    \begin{enumerate}
        \item
              $\ker D$ 称为 $D$ 的\emph{常数环}；

        \item
              若对任意 $a\in A$
              存在 $k$ 使
              \[
                  D^k(a)=0,
              \]
              则称 $D$ 为\emph{局部幂零导子}；

        \item
              若 $D$ 局部幂零，
              定义映射
              \[
                  \phi(a)=
                  \sum_{n\ge0}
                  \frac{D^n(a)}{n!}x^n,
              \]
              则 $\phi:A\to A[x]$
              为环同态，
              并满足  $\phi\circ D=\partial_x\circ\phi$
    \end{enumerate}
\end{definition}
\begin{example}
    设 $A=k[x_1,\dots,x_n]$，则 $A$ 上的导子 $D_i=\frac{\partial}{\partial x_i}$，$i=1,\dots,n$ 是局部幂零的。只关于$x_2,x_3$的多项式$x_2+x_3^5$自然是 $\partial_{x_1}$ 的常数，但不是 $\partial_{x_2}$ 的常数。
\end{example}
\begin{proposition}[见\cite{wright1981jacobian}]
    \label{pro:wright1981jacobia}
    设
    $D_1,\dots,D_n$
    为两两可交换的局部幂零导子，
    且存在
    $t_1,\dots,t_n\in A$
    满足
    \[
        D_i(t_j)=\delta_{ij}。
    \]
    则：

    \begin{enumerate}
        \item  $A=R[t_1,\dots,t_n],$
              其中 $R$ 为$D_1,\dots,D_n$的常数环；

        \item
              $t_i$ 在 $R$ 上代数无关；

        \item
              $ D_i=\frac{\partial}{\partial t_i}$
    \end{enumerate}
\end{proposition}

为了利用\ref{pro:wright1981jacobia}证明Jacobian猜想，我们需要定义适当的导子应用在$f_1,\dots,f_n$上，所需要的导子由下述引理给出。

\begin{lemma}
    $k$-代数$k[x_1,\dots,x_n,\Delta^{-1}]$上的线性映射
    \begin{equation*}
        D_i(\varphi)
        =
        \Delta^{-1}
        \det
        J(f_1,\dots,f_{i-1},\varphi,f_{i+1},\dots,f_n)
    \end{equation*}
    是两两可交换的导子，且满足 $D_i(f_j)=\delta_{ij}$.
\end{lemma}

\begin{theorem}
    若 域$k$上的Weyl 代数 $A_n$ 的每个代数自同态均为自同构，
    则 $k$上的Jacobian 猜想成立。
    \begin{proof}
        只需证明$D_i$作为$k[x_1,\dots,x_n]$上的导子是局部幂零的（注意$k[x_1,\dots,x_n]$上的导子一定含在$A_n$中），然后利用\ref{pro:wright1981jacobia}即可得证。
        存在同态$\psi : A_n \to A_n$使得
        \begin{equation*}
            \psi(x_i)=f_i, \quad \psi(\partial_{x_i})=D_i, \quad i=1,\dots,n.
        \end{equation*}
        根据假设，这$A_n$间的同构。对任意的$D_i\in A_n$和多项式，存在$N$使得
        \begin{equation*}
            \operatorname{ad}_{D_i}^N(x)
            =
            \psi^{-1}\left(\operatorname{ad}_{\psi(D_i)}^N(\psi(x))\right)
            =
            \psi^{-1}\left(\operatorname{ad}_{\partial_{x_i}}^N(\psi(x))\right)
            =
            0
        \end{equation*}
        因为$\partial_{x_i}$是局部幂零的。
    \end{proof}
\end{theorem}



\subsection{$A_n$-模}
$A_n$-模是指 $A_n$ 上的代数模，可以是左模或右模。一般左$A_n$模是我们所需研究的函数空间（即微分算子作用在函数上），而右$A_n$一般是微分形式、分布等对象的模（即微分算子作用在函数的对偶上）。参见\ref{exm:k[x,f^{-1}]}和\cite{bernshtein1972analytic}.
\begin{example}
    \label{exm:k[x]}
    左 $A_n$-模 $k[X]$ 是不可约挠模，进而是一个循环模，由 单位元$1$ 生成。
    且
    \[
        {}_{A_n}k[X]\cong
        A_n/\sum_{i=1}^n A_n\partial_i.
    \]
\end{example}

\begin{proposition}
    商模
    $A_n/\sum_i A_n x_i$
    同构于 $k[\partial_1,\dots,\partial_n]$,
    其中 $\partial_i$ 作用为乘法，
    $x_i$ 的作用为 $-\partial/\partial_i$.
\end{proposition}

\begin{theorem}
    对任意正整数 $r$,
    由
    \[
        \sigma_r(x_i)=x_i,\quad \sigma_r(\partial_i)=\partial_i-x_i^r
    \]
    定义的扭曲模 $k[X]_{\sigma_r}$
    构成一族两两不同构的不可约 $A_n$-模。
    \begin{proof}
        由\ref{thm:A}可得，$\sigma_r$ 定义了一个 $A_n$自同态，因为
        \begin{equation*}
            [\sigma_r(\partial_i),\sigma_r(\partial_j)]
            =[\partial_i-x_i^r,\partial_j-x_j^r]
            =0,
        \end{equation*}
        \begin{equation*}
            [\sigma_r(x_i),\sigma_r(x_j)]
            =[x_i,x_j]
            =0,
        \end{equation*}
        \begin{equation*}
            [\sigma_r(\partial_i),\sigma_r(x_j)]
            =[\partial_i-x_i^r,x_j]
            =\delta_{ij}.
        \end{equation*}
        同时容易看出$\sigma_r$满，结合\ref{thm:simple}，$\sigma_r$是一个自同构。

        由此 $k[X]_{\sigma_r}$ 是良定义的，且由\ref{pro:twisting}可知$k[X]_{\sigma_r}$是一个不可约模，可被$1$生成。
        假设$\phi:k[X]_{\sigma_r}\to k[X]_{\sigma_s}$是一个$A_n$-模同构，则$\phi$完全由$\phi(1)\neq 0$决定。考虑等式
        \begin{equation*}
            -x_i^r\phi(1)
            =
            \phi(-x_i^r)
            =
            \phi(\partial_i\cdot 1)
            =
            (\partial_i-x^s)\phi(1)
        \end{equation*}
        整理得到
        \begin{equation*}
            \partial_i\phi(1)=(x_i^s-x_i^r)\phi(1).
        \end{equation*}
        比较两边的次数，得到矛盾。
    \end{proof}

\end{theorem}


\begin{example}
    对$\mathbb{R}^n$上的光滑函数环$C^\infty(\mathbb{R}^n)$，定义$\partial_i$的作用为$\frac{\partial}{\partial x_i}$，则$C^\infty(\mathbb{R}^n)$是一个左$A_n$-模。而对全体$n$阶微分形式$\Omega^n$，定义$A_n$的右作用为
    \begin{equation*}
        \begin{aligned}
            \omega \cdot x_i        & =  x_i \omega,                    \\
            \omega \cdot \partial_i & =  \mathcal{L}_{\partial_i}\omega
        \end{aligned}
    \end{equation*}
    其中$\mathcal{L}_{\partial_i}$是向量场$\partial_i$对应的Lie导数，则$\Omega^n$是一个右$A_n$-模。
\end{example}
\begin{remark}
    $\left(\mathbb{R}^n,\mathcal{O}_{\mathbb{R}^n}\right)$上的开集$U$对应的微分算子层为$C^\infty(U)\left\langle \partial_1,\ldots,\partial_n \right\rangle$
\end{remark}

\begin{example}
    \label{exm:k[x,f^{-1}]}
    设非零多项式$f\in k[x]$，则在$f$处局部化所得到的模$k[x,f^{-1}]$也是一个左$A_n$-模，其中$\partial_i$的作用为
    \begin{equation*}
        \partial_i\cdot \frac{g}{f^k}=\frac{\partial g}{\partial x_i}\cdot \frac{1}{f^k}-k\cdot \frac{g}{f^{k+1}}\cdot \frac{\partial f}{\partial x_i}.
    \end{equation*}
    $k \in \mathbb{Z}_{\geq 0},\ g \in k[x].$
\end{example}

\begin{example}
    设$U$是$\mathbb{C}$中的一个开集，$\mathcal{H}(U)$表示$U$上所有全纯函数构成的集合，则$\mathcal{H}(U)$是一个左$A_1(\mathbb{C})$-模，$\partial$的作用为$\frac{d}{dx}$。其中$h(z)=e^{e^z}\in \mathcal{H}(U)$不是一个一个挠元素，且$\mathcal{H}(U)$不是循环的：对任意的$p\in \mathcal{H}(U)$, $e^p\notin \mathcal{H}(U)\cdot h$.
    对比 \ref{exm:k[x]}中的例子。
\end{example}



\newpage
